\documentclass[11pt]{article}
\usepackage[T1]{fontenc}
\usepackage[utf8]{inputenc}
\usepackage{lmodern,amsmath,amssymb,amsthm,booktabs}
\usepackage[margin=1in]{geometry}
\usepackage[colorlinks=true,linkcolor=blue,urlcolor=blue,citecolor=blue]{hyperref}
\hypersetup{pdftitle={Paths versus proper double comets for the Randic index times algebraic connectivity},pdfauthor={Jitendra Prajapati}}
\newtheorem{theorem}{Theorem}
\newtheorem{lemma}{Lemma}
\newcommand{\DC}{\mathrm{DC}}
\newcommand{\ra}{\rho}
\title{Paths versus proper double comets for the\protect\\ Randi\'c index times algebraic connectivity}
\author{Jitendra Prajapati\\[3pt]\small Independent\\
\small\href{https://orcid.org/0009-0008-7493-0311}{ORCID: 0009-0008-7493-0311}}
\date{November 5, 2026}
\begin{document}
\maketitle
\begin{abstract}
For each $n\in\{10,11,12,13\}$, we certify that the path $P_n$ has a strictly
smaller product of Randi\'c index and algebraic connectivity than every double
comet whose two end stars each have at least two leaves. At order $14$, the
balanced double comet with two leaves at each end has a smaller product than
the path. Two separate exact-arithmetic programs verify these finite
comparisons, one by spectral intervals and the other by integer determinant
signs. The comparisons contradict the product clause of the 2007
Aouchiche--Hansen--Zheng conjecture under the convention that double comets
exclude paths. The printed definition does not explicitly impose this
restriction; the inclusive convention is therefore distinguished.
\end{abstract}

\section{The product and the source conjecture}
All graphs are finite, simple and connected. For a graph $G$ with vertex
degrees $d_v$, let
\[
 R(G)=\sum_{uv\in E(G)}\frac{1}{\sqrt{d_ud_v}},\qquad
 a(G)=\lambda_2(D-A),\qquad \ra(G)=R(G)a(G),
\]
where the Laplacian eigenvalues are ordered increasingly.
Aouchiche, Hansen and Zheng~\cite[p.~98, Conjecture~4]{AHZ} proposed that
the minimum of $\ra$ among connected graphs of order $n$ is attained by a
path for $n\le9$, and by a balanced double comet for $n\ge10$.
Li, Shi and Wang~\cite[p.~845, Conjecture~1.2]{LSW} repeated the product
clause. Their results on the product impose additional diameter or
minimum-degree hypotheses.

Let $\DC(k_1,k_2,p)$ be the tree formed from two centres joined by a path
with $p$ internal vertices, with $k_1$ pendant leaves at the first centre
and $k_2$ at the second. Thus its order is $k_1+k_2+p+2$.
Throughout the theorem, a \emph{proper double comet} means
$k_1,k_2\ge2$ and $p\ge0$. Balanced double comets with equal star sizes
are a subfamily.

The parenthetical definitions in~\cite{AHZ,LSW} specify equal-sized stars
joined by a path, but no lower bound on their size. If one-leaf stars are
allowed, $\DC(1,1,n-4)$ is the path itself. The split into two ranges of
$n$ does not logically exclude this interpretation. A later convention
in Aouchiche, Hansen, Hertz and Sellal~\cite[Section~2]{AHHS} restricts
balanced double comets to diameter at most $n-3$, which excludes paths,
but does not determine the intended scope of the earlier statement.

\section{A certified finite comparison}
\begin{theorem}\label{thm:main}
For $n\in\{10,11,12,13\}$ and all integers $k_1,k_2\ge2$, $p\ge0$ with
$k_1+k_2+p+2=n$,
\[
 \ra(P_n)<c_n<\ra(\DC(k_1,k_2,p)),
\]
where the rational constants $c_n$ are given in Table~\ref{tab:bounds}.
Moreover,
\[
 \ra(\DC(2,2,8))<\frac{173}{500}<\ra(P_{14}).
\]
\end{theorem}

\begin{table}[ht]
\centering
\begin{tabular}{rrrr}
\toprule
$n$ & $c_n$ & Unordered classes & Ordered triples\\
\midrule
10 & $49/100$ & 9 & 15\\
11 & $11/25$ & 12 & 21\\
12 & $81/200$ & 16 & 28\\
13 & $373/1000$ & 20 & 36\\
\bottomrule
\end{tabular}
\caption{Exact separating constants. Unordered means $k_1\le k_2$;
interchanging the two centres gives an isomorphic graph.}\label{tab:bounds}
\end{table}

\begin{lemma}\label{lem:sign}
Let $L$ be the Laplacian of a connected graph of order $n$, and let $J$
be the all-ones matrix. For $0<t<n$, positive leading principal minors
of $L+J-tI$ imply $a(G)>t$. A negative determinant of $L+J-tI$ implies
$a(G)<t$.
\end{lemma}
\begin{proof}
The all-ones vector is a zero eigenvector of $L$. On its orthogonal
complement, $J$ vanishes. The eigenvalues of $L+J-tI$ are consequently
$n-t,\lambda_2(L)-t,\ldots,\lambda_n(L)-t$.
Sylvester's criterion proves the first assertion. For the second, a
negative determinant forces a negative eigenvalue, and $n-t$ is positive.
\end{proof}

\begin{proof}[Computer-assisted proof of Theorem~\ref{thm:main}]
The program \texttt{verify\_principal\_minors.py} in the accompanying
repository~\cite{artifact} implements the following finite certificate
using integer and rational arithmetic only.

For each edge, put $m=d_ud_v$, $s=10^{12}$, and $q=\lfloor s\sqrt m\rfloor$,
computed as an integer square root of $ms^2$. Then
\[
 \frac{s}{q+1}\le\frac1{\sqrt m}\le\frac{s}{q}.
\]
The program also checks these inequalities by squaring the positive
rational endpoints. Summing gives positive bounds $R_-\le R(G)\le R_+$.
To prove $\ra(G)>c$, set $t=c/R_-$ and check the first criterion of
Lemma~\ref{lem:sign}. To prove $\ra(G)<c$, set $t=c/R_+$ and check the
second. Both uses explicitly check $0<t<n$.

For $t=a/b$ with $b>0$, the integer matrix
\[
 M=b(L+J)-aI
\]
has the same principal-minor signs as $L+J-tI$. Each required determinant
is computed by fraction-free Bareiss elimination, with row swaps and
exact-division checks. For each order in Table~\ref{tab:bounds}, the
program constructs $P_n$ and enumerates all $k_1\le k_2$ with
$k_1,k_2\ge2$ and $k_1+k_2+2\le n$, setting $p=n-k_1-k_2-2$.
Every path gives the required negative determinant at $c_n/R_+$, and
every proper double comet gives positive leading principal minors at
$c_n/R_-$. At order $14$ the path and $\DC(2,2,8)$ satisfy the opposite
comparisons at $173/500$.

These are 63 product comparisons: 57 proper double comets, four paths,
and two order-14 comparisons. The program checks the enumeration counts,
graph simplicity, connectedness, degree multisets and edge counts, and
records each rational spectral threshold. Every comparison succeeds.
Its checks use explicit exceptions and remain active in optimized Python.
\end{proof}

A second program, \texttt{verify\_theorem.py}, independently bounds each
algebraic connectivity in two ways: rational inertia counts for $L-tI$,
and Sturm isolation for the characteristic polynomial after removing its
simple zero factor. It computes characteristic polynomials by the exact
Faddeev--LeVerrier recurrence and checks Cayley--Hamilton. For paths it
also checks
\[
 R(P_n)=\frac{n-3}{2}+\sqrt2,\qquad
 a(P_n)=2-2\cos(\pi/n)
\]
using rational bounds for the radicals, $\pi$ and cosine. Optional SymPy
characteristic polynomials provide a further implementation check.
Floating point supplies trial endpoints and a noncertifying closed-form
sanity check; all final product comparisons use exact rational bounds.

The interval program checks both orders of the star sizes. It also
certifies that $\DC(2,2,n-6)$ is the unique smallest proper-double-comet
class for each $n=10,\ldots,14$, by strict separation of its upper
product bound from every other class's lower bound. At order $14$ all
other proper-double-comet classes are still above the path.

\section{Interpretation and supplementary computation}
Under the proper convention, Theorem~\ref{thm:main} excludes every
balanced double comet from being a global minimizer at orders $10$--$13$:
the path is a connected graph of the same order with a smaller value.
Proving that the path is globally optimal is unnecessary for this
conclusion. Under the inclusive convention, the theorem does not refute
the conjecture, because the path belongs to the named family.
The order-$14$ comparison does not establish a threshold valid for all
larger orders.

Separate saved floating-point scans report $P_{10}$ as the minimizer
among 11,716,571 connected graphs of order $10$. Tree scans through order
$22$ report paths for orders $10$--$13$, $\DC(2,2,n-6)$ for $14$--$17$,
$\DC(3,3,n-8)$ for $18$--$21$, and $\DC(4,4,12)$ for $22$.
These are numerical observations, not exact global optimality
certificates. The repository includes a replay of 40 retained winners
and runners-up, not a rerun of the full enumeration. The interrupted
connected-graph scan at order $11$ is excluded.

Gebendorfer's tadpole construction~\cite{Geb} concerns the quotient
$R/a$, a separate clause. It should not be conflated with these product
comparisons. The accompanying literature record identifies the sources
searched; no earlier instance of this precise finite comparison was
located, and no exhaustive historical-priority claim is made.

\section*{Reproducibility}
The exact verifiers require Python 3.10 or later and no external packages:
\begin{verbatim}
python3 -I certificate/verify_theorem.py --no-sympy
python3 -I -O certificate/verify_principal_minors.py
\end{verbatim}
The first command ends with \texttt{ALL CERTIFIED: True}; the second
outputs a JSON record with \texttt{all\_certified} equal to \texttt{true}.
Both return a nonzero exit status on a failed check. Source code, replay
outputs, review provenance, and this note are available in~\cite{artifact}.
The note is licensed under Creative Commons Attribution 4.0 International.

\begin{thebibliography}{9}
\bibitem{AHZ}
M.~Aouchiche, P.~Hansen, and M.~Zheng,
Variable Neighborhood Search for Extremal Graphs. 19. Further Conjectures
and Results about the Randi\'c Index,
\emph{MATCH Commun. Math. Comput. Chem.} \textbf{58} (2007), 83--102.
\href{https://match.pmf.kg.ac.rs/electronic_versions/Match58/n1/match58n1_83-102.pdf}{Publisher PDF}.
\bibitem{LSW}
X.~Li, Y.~Shi, and L.~Wang,
On a Relation Between Randi\'c Index and Algebraic Connectivity,
\emph{MATCH Commun. Math. Comput. Chem.} \textbf{68} (2012), 843--849.
\href{https://match.pmf.kg.ac.rs/electronic_versions/Match68/n3/match68n3_843-849.pdf}{Publisher PDF}.
\bibitem{AHHS}
M.~Aouchiche, P.~Hansen, A.~Hertz, and C.~Sellal,
\emph{Bounds on differences between some graph theoretic invariants},
Les Cahiers du GERAD G--2015--31, April 2015.
\href{https://www.gerad.ca/fr/papers/G-2015-31.pdf}{Report}.
\bibitem{Geb}
J.~J.~Gebendorfer,
\emph{An infinite tadpole counterexample family to the Randi\'c
index--algebraic connectivity conjecture}, preprint, August 6, 2026.
\href{https://doi.org/10.5281/zenodo.21827918}{doi:10.5281/zenodo.21827918}.
\bibitem{artifact}
J.~Prajapati, \emph{Paths versus proper double comets for Randi\'c index
times algebraic connectivity}, computational artifact, version 1.0.0,
November 5, 2026.
\url{https://github.com/infinityscroll/agx-randic-connectivity-certificate}.
\end{thebibliography}
\end{document}
